% GENERATED FILE. Edit canonical lessons, then run npm run generate:books.
% canonical-source-hash: fnv1a64:7b233b3acdbdbf60
% canonical-lessons: ZH-C54-xuexi, ZH-C54-xie, ZH-C54-du, ZH-C54-wen, ZH-C54-ke

\chapter{Learning}
\label{ch:zh-learning}
\hlchaptermodality{54}

\begin{chapteropening}
\textbf{By the end of this chapter:} \emph{I can say what I study, read and ask about in class.}

Complete the last lesson of chapter 54: I can say what I study, read and ask about in class.
\end{chapteropening}

\section[xuéxí]{xuéxí — to study; to learn}

\label{lesson:ZH-C54-xuexi}



\begin{quote}
Before the new one: say the Mandarin for to understand, then the Mandarin for slow.
\end{quote}

\subsection*{You'll want to know: xuéxí}
\emph{xuéxí} — \textquotedblleft{}to study; to learn\textquotedblright{}.

\begin{quote}
\emph{wǒ xuéxí hànyǔ} — I'm learning Chinese.
\end{quote}

Its characters come later; for now it is a word for the ear.

\begin{cousinweb}
\emph{xuéxí} is built from \zh{学} 'to learn' and xí 'to practise', both inherited from Middle Chinese, \emph{haewk} and \emph{zip} (Guangyun, in Baxter's transcription); the dictionary glosses \emph{zip} with the first.
\end{cousinweb}

\begin{grammarlens}[title={What you've built}]
One more word for asking politely.
\end{grammarlens}

\subsection*{Your turn}
\begin{itemize}
\raggedright
  \item \emph{Say it:} \emph{xuéxí}
  \item \emph{Say it:} \emph{xuéxí}, once more
  \item \emph{Say it:} \emph{wǒ xuéxí hànyǔ} — I'm learning Chinese.
  \item \emph{Recall:} say the Mandarin for to understand, then the Mandarin for slow, then say \emph{xuéxí} again
\end{itemize}

\begin{tcolorbox}[breakable,colback=teal!4,colframe=teal!35!black,title={Before you move on}]
What does xuéxí mean? (\textquotedblleft{}To study; to learn\textquotedblright{}.) Say it once more.
\end{tcolorbox}
\section[xiě]{xiě — to write}

\label{lesson:ZH-C54-xie}



\begin{quote}
Before the new one: say the Mandarin for slow, then the Mandarin for to learn.
\end{quote}

\subsection*{You'll want to know: xiě}
\emph{xiě} — \textquotedblleft{}to write\textquotedblright{}.

\begin{quote}
\emph{wǒ xiě hànzì} — I'm practising characters.
\end{quote}

Its characters come later; for now it is a word for the ear.

\begin{cousinweb}
\emph{xiě} is inherited from Middle Chinese \emph{sjaeX} (Guangyun, in Baxter's transcription), glossed 'to remove' and 'to copy out a book'.
\end{cousinweb}

\begin{grammarlens}[title={What you've built}]
One more word for asking politely.
\end{grammarlens}

\subsection*{Your turn}
\begin{itemize}
\raggedright
  \item \emph{Say it:} \emph{xiě}
  \item \emph{Say it:} \emph{xiě}, once more
  \item \emph{Say it:} \emph{wǒ xiě hànzì} — I'm practising characters.
  \item \emph{Recall:} say the Mandarin for slow, then the Mandarin for to learn, then say \emph{xiě} again
\end{itemize}

\begin{tcolorbox}[breakable,colback=teal!4,colframe=teal!35!black,title={Before you move on}]
What does xiě mean? (\textquotedblleft{}To write\textquotedblright{}.) Say it once more.
\end{tcolorbox}
\section[dú]{dú — to read aloud; to study}

\label{lesson:ZH-C54-du}



\begin{quote}
Before the new one: say the Mandarin for to learn, then the Mandarin for to write.
\end{quote}

\subsection*{You'll want to know: dú}
\emph{dú} — \textquotedblleft{}to read aloud; to study\textquotedblright{}.

\begin{quote}
\emph{wǒ xǐhuan dú shū} — I love reading.
\end{quote}

Its characters come later; for now it is a word for the ear.

\begin{cousinweb}
\emph{dú} is inherited from Middle Chinese \emph{duwk} (Guangyun, in Baxter's transcription), glossed 'to recite'.
\end{cousinweb}

\begin{grammarlens}[title={What you've built}]
One more word for asking politely.
\end{grammarlens}

\subsection*{Your turn}
\begin{itemize}
\raggedright
  \item \emph{Say it:} \emph{dú}
  \item \emph{Say it:} \emph{dú}, once more
  \item \emph{Say it:} \emph{wǒ xǐhuan dú shū} — I love reading.
  \item \emph{Recall:} say the Mandarin for to learn, then the Mandarin for to write, then say \emph{dú} again
\end{itemize}

\begin{tcolorbox}[breakable,colback=teal!4,colframe=teal!35!black,title={Before you move on}]
What does dú mean? (\textquotedblleft{}To read aloud; to study\textquotedblright{}.) Say it once more.
\end{tcolorbox}
\section[wèn]{wèn — to ask}

\label{lesson:ZH-C54-wen}



\begin{quote}
Before the new one: say the Mandarin for to write, then the Mandarin for to read aloud.
\end{quote}

\subsection*{You'll want to know: wèn}
\emph{wèn} — \textquotedblleft{}to ask\textquotedblright{}.

\begin{quote}
\emph{wǒ wèn lǎoshī} — I'll ask the teacher.
\end{quote}

Its characters come later; for now it is a word for the ear.

\begin{cousinweb}
\emph{wèn} is inherited from Middle Chinese \emph{mjunH} (Guangyun, in Baxter's transcription), glossed 'to inquire'.
\end{cousinweb}

\begin{grammarlens}[title={What you've built}]
One more word for asking politely.
\end{grammarlens}

\subsection*{Your turn}
\begin{itemize}
\raggedright
  \item \emph{Say it:} \emph{wèn}
  \item \emph{Say it:} \emph{wèn}, once more
  \item \emph{Say it:} \emph{wǒ wèn lǎoshī} — I'll ask the teacher.
  \item \emph{Recall:} say the Mandarin for to write, then the Mandarin for to read aloud, then say \emph{wèn} again
\end{itemize}

\begin{tcolorbox}[breakable,colback=teal!4,colframe=teal!35!black,title={Before you move on}]
What does wèn mean? (\textquotedblleft{}To ask\textquotedblright{}.) Say it once more.
\end{tcolorbox}
\section[kè]{kè — lesson; class}

\label{lesson:ZH-C54-ke}



\begin{quote}
Before the new one: say the Mandarin for to read aloud, then the Mandarin for to ask.
\end{quote}

\subsection*{You'll want to know: kè}
\emph{kè} — \textquotedblleft{}lesson; class\textquotedblright{}.

\begin{quote}
\emph{jīntiān yǒu hànyǔ kè} — There's a Chinese class today.
\end{quote}

Its characters come later; for now it is a word for the ear.

\begin{cousinweb}
\emph{kè} is inherited from Middle Chinese \emph{khwaH} (Guangyun, in Baxter's transcription), glossed 'a tax' and 'a test'. The sense 'lesson' is not given there, so how it got it is not settled.
\end{cousinweb}

\begin{grammarlens}[title={What you've built}]
One more word for asking politely.
\end{grammarlens}

\subsection*{Your turn}
\begin{itemize}
\raggedright
  \item \emph{Say it:} \emph{kè}
  \item \emph{Say it:} \emph{kè}, once more
  \item \emph{Say it:} \emph{jīntiān yǒu hànyǔ kè} — There's a Chinese class today.
  \item \emph{Recall:} say the Mandarin for to read aloud, then the Mandarin for to ask, then say \emph{kè} again
\end{itemize}

\begin{tcolorbox}[breakable,colback=teal!4,colframe=teal!35!black,title={Before you move on}]
What does kè mean? (\textquotedblleft{}Lesson; class\textquotedblright{}.) Say it once more.
\end{tcolorbox}
